\SeventhArticle{8}
Um nun das Integral $\int Y\,\mathrm{d}t$ zu bestimmen, betrachten wir die einzelnen Teile von $Y$. Wir bezeichnen das Produkt
\[
(t-a)(t-a')(t-a'')(t-a''')\ldotp\ldotp\ldotp\ldotp(t-a^{(n)})
\]
mit $T$. Durch Ausmultiplizieren dieses Produkts ergibt sich
\[
T=t^{n+1}+\alpha t^n+\alpha't^{n-1}+\alpha''t^{n-2}+\text{usw.}+\alpha^{(n)}
\]
Der Zähler des Bruchs, mit dem im ersten Teil von $Y$ die Größe $A$ multipliziert ist, lautet $=\frac{T}{t-a}$; die Zähler in den folgenden Teilen sind entsprechend $\frac{T}{t-a'},\frac{T}{t-a''},\frac{T}{t-a'''}$ usw. Die Nenner sind dagegen nichts anderes als die Werte dieser Zähler, wenn jeweils $t=a$, $t=a'$, $t=a''$, $t=a'''$ usw. gesetzt wird. Wir bezeichnen diese Nenner entsprechend mit $M,M',M'',M'''$ usw., sodass
\[
\begin{aligned}&Y=\frac{AT}{M(t-a)}+\frac{A'T}{M'(t-a')}+\frac{A''T}{M''(t-a'')}\\ &+\text{usw.}+\frac{A^{(n)}T}{M^{(n)}(t-a^{(n)})}\end{aligned}
\]
Da $T=0$ für $t=a$ ist, haben wir die identische Gleichung
\[
a^{n+1}+\alpha a^n+\alpha'a^{n-1}+\alpha''a^{n-2}+\text{usw.}+\alpha^{(n)}=0
\]
und daher
\[
\begin{aligned}
T={}&t^{n+1}-a^{n+1}+\alpha(t^n-a^n)+\alpha'(t^{n-1}-a^{n-1})
\\ &+\alpha''(t^{n-2}-a^{n-2})+\text{usw.}\\
&+\alpha^{(n-1)}(t-a)
\end{aligned}
\]
Division durch $t-a$ ergibt nun
\[
\begin{aligned}
\frac{T}{t-a}={}&t^n+at^{n-1}+aat^{n-2}+a^3t^{n-3}+\text{usw.}+a^n\\[4pt]
&+\alpha t^{n-1}+\alpha at^{(n-2)}+\alpha aat^{n-3}+\text{usw.}+\alpha a^{n-1}\\[4pt]
&\hspace{1.01em}+\alpha't^{n-2}+\alpha'at^{n-3}+\text{usw.}+\alpha'a^{n-2}\\[4pt]
&\hspace{2.01em}+\alpha''t^{n-3}+\text{usw.}+\alpha''a^{n-3}\\[4pt]
&\hspace{3.01em}+\text{usw. usw.}\\[4pt]
&\hspace{4.01em}+\alpha^{(n-1)}
\end{aligned}
\]
\SourceLeaf{183}Der Wert dieser Funktion für $t=a$ ergibt sich zu
\[
\begin{aligned}&=(n+1)a^n+n\alpha a^{n-1}+(n-1)\alpha'a^{n-2}\\ &+(n-2)\alpha''a^{n-3}
+\text{usw.}+\alpha^{(n-1)}\end{aligned}
\]
Somit ist $M$ gleich dem Wert von $\frac{\mathrm{d}T}{\mathrm{d}t}$ für $t=a$, wie auch aus einer anderen Betrachtung bekannt ist. Entsprechend sind $M',M'',M'''$ usw. die Werte von $\frac{\mathrm{d}T}{\mathrm{d}t}$ für $t=a'$, $t=a''$, $t=a'''$ usw.

Ferner finden wir für das Integral $\int\frac{T\,\mathrm{d}t}{t-a}$ von $t=0$ bis $t=1$ den Wert
\[
\begin{aligned}
={}&\frac{1}{n+1}+\frac{a}{n}+\frac{aa}{n-1}+\frac{a^3}{n-2}+\text{usw.}+a^n\\[4pt]
&+\frac{\alpha}{n}+\frac{\alpha a}{n-1}+\frac{\alpha aa}{n-2}+\text{usw.}+\alpha a^{n-1}\\[4pt]
&\hspace{5.01em}+\frac{\alpha'}{n-1}+\frac{\alpha'a}{n-2}+\text{usw.}+\alpha'a^{n-2}\\[4pt]
&\hspace{6.01em}+\frac{\alpha''}{n-2}+\text{usw.}+\alpha''a^{n-3}\\[4pt]
&\hspace{7.01em}+\text{usw. usw.}\\[4pt]
&\hspace{8.01em}+\alpha^{(n-1)}
\end{aligned}
\]
Diese Glieder ordnen wir folgendermaßen an:
\[
\begin{aligned}
&a^n+\alpha a^{n-1}+\alpha'a^{n-2}+\alpha''a^{n-3}+\text{usw.}+\alpha^{(n-1)}\\[4pt]
&\quad+\tfrac12(a^{n-1}+\alpha a^{n-2}+\alpha'a^{n-3}+\text{usw.}+\alpha^{(n-2)})\\[4pt]
&\quad+\tfrac13(a^{n-2}+\alpha a^{n-3}+\alpha'a^{n-4}+\text{usw.}+\alpha^{(n-3)})\\[4pt]
&\quad+\tfrac14(a^{n-3}+\alpha a^{n-4}+\alpha'a^{n-5}+\text{usw.}+\alpha^{(n-4)})\\[4pt]
&\quad+\text{usw.}\\[4pt]
&\quad+\frac{1}{n-1}(aa+\alpha a+\alpha')\\[4pt]
&\quad+\frac{1}{n}(a+\alpha)\\[4pt]
&\quad+\frac{1}{n+1}
\end{aligned}
\]
Offenbar ergibt sich aber dieselbe Größe, wenn in dem Produkt aus der Multiplikation der Funktion $T$ mit der unendlichen Reihe
\[
t^{-1}+\tfrac12t^{-2}+\tfrac13t^{-3}+\tfrac14t^{-4}\ \text{usw.}
\]
nach Weglassen aller Glieder, die Potenzen von $t$ mit negativen Exponenten enthalten (oder kürzer: in dem Teil des Produkts, der ein Polynom in $t$ ist), $t$ durch $a$ ersetzt wird. Wir setzen daher\HistoricalNote{19}
\[
T(t^{-1}+\tfrac12t^{-2}+\tfrac13t^{-3}+\tfrac14t^{-4}+\text{usw.})=T'+T''
\]
\SourceLeaf{184}Dabei sei $T'$ das in diesem Produkt enthaltene Polynom in $t$, während $T''$ den anderen Teil bezeichnet, nämlich eine unendliche Reihe nach absteigenden Potenzen. Dann ist der Wert des Integrals $\int\frac{T\,\mathrm{d}t}{t-a}$ von $t=0$ bis $t=1$ gleich dem Wert der Funktion $T'$ für $t=a$. Bezeichnen wir also die Werte der Funktion
\[
\frac{T'}{\left(\dfrac{\mathrm{d}T}{\mathrm{d}t}\right)}
\]
für $t=a$, $t=a'$, $t=a''$, $t=a'''$ usw. bis $t=a^{(n)}$ entsprechend mit $R,R',R'',R'''\ldotp\ldotp\ldotp\ldotp R^{(n)}$, so ergibt sich für das Integral $\int Y\,\mathrm{d}t$ von $t=0$ bis $t=1$
\[
=RA+R'A'+R''A''+\text{usw.}+R^{(n)}A^{(n)}
\]
was, mit $\Delta$ multipliziert, entweder den wahren oder den angenäherten Wert des Integrals $\int y\,\mathrm{d}x$ von $x=g$ bis $x=g+\Delta$ liefert.

\SeventhArticle{9}
Diese Operationen lassen sich etwas leichter ausführen, wenn anstelle der Unbestimmten $t$ eine andere, $u=2t-1$, eingeführt wird. Der Kürze wegen schreiben wir außerdem $b=2a-1$, $b'=2a'-1$, $b''=2a''-1$ usw. Wird in $T$ anstelle von $t$ der Wert $\tfrac12u+\tfrac12$ eingesetzt, so gehe die Funktion in $\frac{U}{2^{n+1}}$ über; es sei also
\[
U=(u-b)(u-b')(u-b'')\ldotp\ldotp\ldotp\ldotp(u-b^n)
\]
Dann ist $\frac{\mathrm{d}T}{\mathrm{d}t}=\frac{1}{2^n}\cdot\frac{\mathrm{d}U}{\mathrm{d}u}$, und somit sind $M,M',M''$ usw. die Werte von $\frac{1}{2^n}\cdot\frac{\mathrm{d}U}{\mathrm{d}u}$, wenn der Reihe nach $u=b$, $u=b'$, $u=b''$ usw. gesetzt wird.

Da die Reihe $t^{-1}+\tfrac12t^{-2}+\tfrac13t^{-3}+\tfrac14t^{-4}+\text{usw.}$ nichts anderes ist als $\log.\frac{1}{1-t^{-1}}=\log\frac{1+u^{-1}}{1-u^{-1}}$, geht sie durch die Substitution $t=\tfrac12u+\tfrac12$ notwendig in $2u^{-1}+\tfrac23u^{-3}+\tfrac25u^{-5}+\tfrac27u^{-7}+\text{usw.}$ über. Setzen wir daher
\[
U(u^{-1}+\tfrac13u^{-3}+\tfrac15u^{-5}+\tfrac17u^{-7}+\text{usw.})=U'+U''
\]
wobei $U'$ das in diesem Produkt enthaltene Polynom in $u$ und $U''$ den anderen Teil, nämlich eine unendliche Reihe nach absteigenden Potenzen, bezeichnet, so ist offenbar
\[
T'+T''=\frac{1}{2^n}(U'+U'')
\]
Aber $T'$ muss als Polynom in $t$ offenbar durch die Substitution $t=\tfrac12u+\tfrac12$ ein Polynom in $u$ ergeben. Dagegen erzeugt $T''$, das nur negative Potenzen von $t$ enthält, durch dieselbe Substitution auch nur negative Potenzen\SourceLeaf{185} von $u$. Somit ist $U'$ nichts anderes als das durch diese Substitution umgeformte $2^nT'$, und entsprechend entsteht $U''$ aus $2^nT''$. Es ist daher gleichgültig, ob wir in $\dfrac{T'}{\left(\dfrac{\mathrm{d}T}{\mathrm{d}t}\right)}$ den Wert $t=a$ einsetzen oder in $\dfrac{U'}{\left(\dfrac{\mathrm{d}U}{\mathrm{d}u}\right)}$ die Gleichung $u=b$ ansetzen. Daraus schließen wir: $R, R', R'', R'''$ usw. sind auch die Werte der Funktion $\dfrac{U'}{\left(\dfrac{\mathrm{d}U}{\mathrm{d}u}\right)}$ für $u=b$. $u=b',\ u=b'',\ u=b'''$ usw.

\SeventhArticle{10}
Bevor wir weitergehen, wollen wir diese Vorschriften an einem Beispiel veranschaulichen. Es sei $n=5$, und wir setzen $a=0,\ a'=\frac15,\ a''=\frac25,\ a'''=\frac35,\ a''''=\frac45,\ a'''''=1$. Daraus ergibt sich
\[
T=t^6-3t^5+\tfrac{17}{5}t^4-\tfrac95t^3+\tfrac{274}{625}tt-\tfrac{24}{625}t
\]
Durch Multiplikation mit $t^{-1}+\frac12t^{-2}+\frac13t^{-3}+\frac14t^{-4}+$ usw. erhalten wir
\[
T'=t^5-\tfrac52t^4+\tfrac{67}{30}t^3-\tfrac{17}{20}tt+\tfrac{913}{7500}t-\tfrac{19}{7500}
\]
Die Werte der Koeffizienten $R, R', R'', R''', R'''', R'''''$ werden daher durch die rationale Funktion
\[
\frac{\begin{gathered}t^5-\frac25t^4+\frac{67}{30}t^3\\-\frac{17}{20}tt+\frac{913}{7500}t-\frac{19}{7500}\end{gathered}}
{\begin{gathered}6t^5-15t^4+\frac{68}{5}t^3\\-\frac{27}{5}tt+\frac{548}{625}t-\frac{24}{625}\end{gathered}}
\]
ausgedrückt, in der für $t$ der Reihe nach die Werte $0,\ \frac15,\ \frac25,\ \frac35,\ \frac45,\ 1$ einzusetzen sind. Etwas kürzer ist die andere Methode, die $b=-1,\ b'=-\frac35,\ b''=-\frac15,\ b'''=\frac15,\ b''''=\frac35,\ b'''''=1$ liefert.
\[
\begin{aligned}
U&=u^6-\tfrac75u^4+\tfrac{259}{625}uu-\tfrac9{625}\\
U'&=u^5-\tfrac{16}{15}u^3+\tfrac{277}{1875}u
\end{aligned}
\]
Daraus ergeben sich $R.\ R',\ R''$ usw. als Werte der rationalen Funktion
\[
\frac{u^4-\frac{16}{15}uu+\frac{277}{1875}}
{6u^4-\frac{28}{5}uu+\frac{518}{625}}
\]
für $u=-1,\ u=-\frac35,\ u=-\frac15$ usw. Beide Methoden liefern dieselben Zahlen, die wir in Artikel 4 aus der Harmonia Mensurarum mitgeteilt haben. In einem Fall wie dem vorliegenden Beispiel, in dem $a, a', a''$ usw. rationale Größen sind, lassen sich die Werte des Nenners $\frac{\mathrm{d}T}{\mathrm{d}t}$ übrigens bequemer in seiner ursprünglichen Gestalt berechnen, nämlich $(a-a')(a-a'').(a-a''')\ldotp\ldotp\ldotp\ldotp(a-a^{(n)})$ für $t=a$, und entsprechend für die übrigen. Dasselbe gilt für den Nenner $\frac{\mathrm{d}U}{\mathrm{d}u}$, der für $u=b$ den Wert $=(b-b')(b-b'')(b-b''')\ldotp\ldotp\ldotp\ldotp(b-b^{(n)})$ annimmt.
